\documentclass[10pt,a4paper]{amsart}
\usepackage[margin=19mm]{geometry}
\usepackage{amsmath,amssymb,mathtools}
\usepackage[scr=rsfs,cal=euler]{mathalfa}
\usepackage{xfrac}
\usepackage[unicode,hidelinks,bookmarks=false]{hyperref}
\newcommand{\ie}{i.e.\ }
\newcommand{\cf}{cf.\ }
\newcommand{\resp}{respectively\ }
\newcommand{\Subsection}{\subsection}
\providecommand{\nobreakdash}{\nobreak\hbox{-}}
\setlength{\emergencystretch}{2em}
\multlinegap0pt
\usepackage{mathtools}
\usepackage{tikz-cd}
\newtheorem{thm}{Theorem}[section]
\newtheorem{intro}{Theorem}
\newtheorem{cor}[thm]{Corollary}
\newtheorem{lem}[thm]{Lemma}
\newtheorem{prop}[thm]{Proposition}
\newtheorem{defn}[thm]{Definition}
\newtheorem{rem}[thm]{Remark}
\newtheorem{construction}[thm]{Construction}
\newtheorem{observation}[thm]{Observation}
\newtheorem{exam}[thm]{Example}
\numberwithin{equation}{section}
\newcommand{\mycomment}[1]{}
\def\dashmapsto{\mapstochar\dashrightarrow}
\newcommand{\norm}[1]{\left\Vert#1\right\Vert}
\newcommand{\abs}[1]{\left\vert#1\right\vert}
\newcommand{\set}[1]{\left\{#1\right\}}
\newcommand{\RR}{\mathbb R}
\newcommand{\QQ}{\mathbb Q}
\newcommand{\NN}{\mathbb N}
\newcommand{\ZZ}{\mathbb Z}
\newcommand{\CC}{\mathbb C}
\newcommand{\PP}{\mathbb P}
\newcommand{\FF}{\mathbb F}
\newcommand{\eps}{\varepsilon}
\newcommand{\lra}{\longrightarrow}
\newcommand{\hra}{\hookrightarrow}
\newcommand{\ra}{\rightarrow}
\newcommand{\la}{\leftarrow}
\newcommand{\ud}{\underline{d}}
\newcommand{\cA}{\mathcal{A}}
\newcommand{\cL}{\mathcal{L}}
\newcommand{\cJ}{\mathcal{J}}
\newcommand{\cK}{\mathcal{K}}
\newcommand{\cH}{\mathcal{H}}
\newcommand{\cU}{\mathcal{U}}
\newcommand{\cT}{\mathcal{T}}
\newcommand{\tC}{\Tilde{C}}
\newcommand{\tX}{\widetilde{X}}
\newcommand{\tY}{\widetilde{Y}}
\newcommand{\tZ}{\widetilde{Z}}
\newcommand{\tE}{\widetilde{E}}
\newcommand{\tP}{\widetilde{P}}
\newcommand{\tT}{\widetilde{T}}
\newcommand{\cC}{\mathcal{C}}
\newcommand{\cE}{\mathcal{E}}
\newcommand{\cF}{\mathcal{F}}
\newcommand{\cG}{\mathcal{G}}
\newcommand{\cM}{\mathcal{M}}
\newcommand{\cN}{\mathcal{N}}
\newcommand{\cV}{\mathcal{V}}
\newcommand{\cO}{\mathcal{O}}
\newcommand{\cD}{\mathcal{D}}
\newcommand{\cP}{\mathcal{P}}
\newcommand{\cR}{\mathcal{R}}
\newcommand{\cRH}{\mathcal{RH}}
\newcommand{\cS}{\mathcal{S}}
\newcommand{\cQ}{\mathcal{Q}}
\newcommand{\cZ}{\mathcal{Z}}
\newcommand{\cX}{\mathcal{X}}
\newcommand{\cY}{\mathcal{Y}}
\newcommand{\cB}{\mathcal{B}}
\newcommand{\ch}{\mathfrak{h}}
\newcommand{\fD}{\mathfrak{D}}
\newcommand{\s}{\sigma}
\newcommand{\st}{\sigma\tau}
\newcommand{\ee}{\mathcal{E}_2(2,2)}
\usepackage{todonotes}
\setlength{\marginparwidth}{0.8in}
\newcommand{\astodo}[1]{\todo[color=blue!40,bordercolor=blue,size=\footnotesize]{\textbf{Anatoli: }#1}}
\DeclareMathOperator{\Aut}{Aut}
\DeclareMathOperator{\id}{id}
\DeclareMathOperator{\Fix}{Fix}
\DeclareMathOperator{\im}{Im}
\DeclareMathOperator{\Nm}{Nm}
\DeclareMathOperator{\Pic}{Pic}
\DeclareMathOperator{\End}{End}
\DeclareMathOperator{\Ker}{Ker}
\DeclareMathOperator{\Hom}{Hom}
\DeclareMathOperator{\Image}{Im}
\DeclareMathOperator{\Spec}{Spec}
\newcommand{\unnumberedfootnote}[1]{%
  \begingroup
  \footnotetext[0]{#1}%
  \endgroup
}
\begin{document}
\title[A note on smooth quotients of Prym varieties]{A note on smooth quotients of Prym varieties}
\author{Anatoli Shatsila}
\date{}
\maketitle
\noindent\emph{Source and licence.} A note on smooth quotients of Prym varieties. Version arXiv:2608.18000v1 (CC BY 4.0). This material is distributed under the Creative Commons Attribution 4.0 International licence, http://creativecommons.org/licenses/by/4.0/, as recorded in the arXiv OAI licence metadata for this exact version and shown on the abstract page https://arxiv.org/abs/2608.18000v1. Full rights notice: licenses/arxiv-2608.18000-CC-BY-4.0.txt.\par\smallskip
\noindent\textit{Editorial scope.} Counted unit: the original \texttt{thm} environment \texttt{thm41} at source lines 347--354 together with its complete proof at lines 356--425; the delivered document keeps source lines 146--426 so that every referenced label and every citation resolves inside it. Native editable source is the paper's own concatenated TeX; the AMS wrapper is editorial formatting only. No publisher class, upstream script or unrelated artwork is redistributed.
\begin{thm}[Theorem \ref{thm31}]\label{thm:order-two}
Let $g\geq4$. Every pseudoreflection on $P(\widetilde C/C)$ of geometric origin has order $2$.
\end{thm}

Moreover, we show that $g \geq 4$ is sharp (see Example \ref{ex:order-four}). The main idea of the proof is a careful analysis of the ramification data of intermediate covers from the properties of pseudoreflections.  

Combining Theorem \ref{thm:order-two} with the list in \cite[Proposition 5.1]{ALN} and the classification in \cite[Theorem 3.5]{ALA} yields the following.

\begin{thm}[Theorem \ref{thm41}]\label{thm:main}
Let $g\geq5$ and let $G\leq \Aut(\widetilde C)$ be a finite group such that every element of $G$ commutes with $\iota$. Let
\[
\rho\colon G\longrightarrow \Aut(P,\Xi)
\]
be the induced representation. Assume that $\rho$ is injective and that $P/\rho(G)$ is smooth. If $G$ is nontrivial, then $$ G\cong \ZZ/2\ZZ \qquad\text{or}\qquad G\cong (\ZZ/2\ZZ)^2.$$
Moreover, if $g > 7$ then $G\cong \ZZ/2\ZZ$. 
Both groups occur for some $g \geq 5$.
\end{thm}

Groups containing elements of order 4 are eliminated by Theorem \ref{thm:order-two}. The situation with 2-groups is more delicate and requires additional arguments involving an analysis of the decomposition of the Prym and its theta divisor. 

\subsection*{Acknowledgements.} 
The author has been supported by the Polish National Science Center project number 2024/53/N/ST1/01634.


\section{Preliminaries}\label{sec:prelim}

\subsection{Prym varieties and geometric automorphisms}

Let $f:\tC\to C$ be an \'{e}tale double cover of a smooth projective complex curve of genus $g$. The Prym variety of the cover $f$ is defined as $$P=P(\widetilde C/C):=\ker\bigl(\operatorname{Nm}_f: J\tC\to JC\bigr)^0.$$
The Prym variety \(P\) carries a natural principal polarization $\Xi$ and we have \(\dim P = d = g-1\). The covering involution $\iota$ acts as $+1$ on $f^*JC$ and as $-1$ on $P$. On holomorphic differentials this gives
$$H^0(\widetilde C,K_{\widetilde C})
=H^0(\widetilde C,K_{\widetilde C})^+
\oplus H^0(\widetilde C,K_{\widetilde C})^-,$$
with
$$ H^0(P,\Omega_P^1)\cong H^0(\widetilde C,K_{\widetilde C})^-.$$
Let $$Z(\iota)=\{\sigma\in\Aut(\widetilde C):\sigma\iota=\iota\sigma\}.$$
Every automorphism $\sigma\in Z(\iota)$ preserves $P$ and induces a polarized automorphism of the Prym. We denote the resulting homomorphism by
\[
\rho\colon Z(\iota)\longrightarrow \Aut(P,\Xi).
\]
An automorphism of $P$ lying in the image of $\rho$ is said to be of geometric origin.

We will use the following result.

\begin{lem}[\cite{ALN}, Lemma 3.1]\label{lem:kernel}
Assume $g\geq5$. Then $|\ker\rho|\leq2$. Moreover, if $\operatorname{id}\neq\mu\in\ker\rho$, then $$g\bigl(\widetilde C/\langle\mu,\iota\rangle\bigr)\leq1.$$
\end{lem}

\subsection{Pseudoreflections}

Let a finite group $H$ act holomorphically on a complex manifold $X$. \textit{A pseudoreflection at $x\in X$} is an element of $H$ whose fixed locus passing through $x$ has pure codimension $1$. If $A$ is an abelian variety and $s$ is an automorphism, then $s$ is a pseudoreflection if and only if its differential on $T_0A$ fixes a hyperplane pointwise (note that in this case we do not need to specify the points $x$). This means that the eigenvalues of $s$ on $T_0A$, and, equivalently, on its dual $H^0(A,\Omega_A^1)$, are
$$1,\ldots,1,\lambda, \qquad \lambda\neq1.$$

The Chevalley-Shephard-Todd theorem says that $X/H$ is smooth if and only if for every $x \in X$ the group $\operatorname{Stab}_H(x)$ is generated by pseudoreflections. Thus, if $H$ is a group of automorphisms of an abelian variety $A$, then smoothness of $A/H$ implies that $H$ is generated by pseudoreflections.

We have the following result.

\begin{prop}[\cite{ALN}, Proposition 3.3]\label{thm:ALN-orders}
For $g\geq3$, every pseudoreflection on a Prym variety of geometric origin has order $2$ or $4$.
\end{prop}

We will eliminate the second possibility when $g\geq4$.

\section{Excluding order four}\label{sec:order4}

We now prove Theorem \ref{thm:order-two}.

\begin{thm}
\label{thm31}
Let $g\geq 4$. Every pseudoreflection on $P(\widetilde C/C)$ of geometric origin has order $2$.
\end{thm}

\begin{proof}
We start with the case $g \geq 5$. By Proposition \ref{thm:ALN-orders}, the only possible orders are $2$ and $4$. Take an order-$4$ pseudoreflection  $\alpha\in\Aut(P,\Xi)$ of geometric origin. Choose $\sigma\in Z(\iota)$ with $\rho(\sigma)=\alpha.$ By Lemma \ref{lem:kernel}, the intersection $\langle\sigma\rangle\cap\ker\rho$ has order at most $2$. Since $\alpha$ has order $4$, it follows that $\sigma$ has order $4$ or $8$.

First, we rule out the case in which $\sigma$ has order $4$. Since $\alpha=\rho(\sigma)$ is an order-$4$ pseudoreflection, its action on $$H^0(P,\Omega_P^1)\cong H^0(\widetilde C,K_{\widetilde C})^-$$
has eigenvalues $$\underbrace{1,\ldots,1}_{d-1},\zeta,
\ \ \  \zeta\in\{i,-i\},$$ where $d=\dim P=g-1\geq4$. In particular, $\rho(\sigma^2)$ has eigenvalues $1^{d-1},-1$, whereas $\rho(\iota)=-\id_P$. Hence $\sigma^2\neq\iota$, and therefore
$$G:=\langle\iota,\sigma\rangle\cong\ZZ/2\ZZ\times\ZZ/4\ZZ.$$
Set $B:=\tC/G$ and consider two intermediate quotient curves
$$C_{\sigma}:=\tC/\langle\sigma\rangle, \ \ \ C_{\iota\sigma}:=\tC/\langle\iota\sigma\rangle.$$

Let us compute their genera. A differential on $C_{\sigma}$ is a $\sigma$-invariant differential on $\widetilde C$. On the $\iota$-invariant part, the $\sigma$-invariants are precisely the $G$-invariants, hence have dimension $g(B)$. The Prym part contributes \(d-1\) as \(\rho(\s)\) is a pseudoreflection. Therefore, $$g(C_{\s})=g(B)+d-1.$$
For $C_{\iota\s}$, a Prym differential fixed by $\iota\sigma$ must satisfy $\sigma^*\omega=-\omega$, but $-1$ is not an eigenvalue of $\rho(\sigma)$. Hence only the $G$-invariant differentials contribute, and we get $g(C_{\iota\sigma}) = g(B)$.

Let $r_{\sigma}$ and $r_{\iota\sigma}$ be the numbers of branch points of the double covers $C_{\sigma} \to B$ and $C_{\iota\sigma}\to B$. The Riemann-Hurwitz formula gives $$r_{\sigma}=2(d-g(B)), \ \ \ r_{\iota\sigma}=2(1-g(B)),$$
so
$$r_{\sigma}-r_{\iota\sigma}=2d-2.$$
Therefore, at least $2d-2$ points of $B$ are ramified in $C_{\sigma}/B$ but unramified in $C_{\iota\sigma}/B$.

Let $b\in B$ be one such point and choose a point $\widetilde b\in \tC$ lying above $b$. We denote by $I_b:=\operatorname{Stab}_G(\widetilde b)
=\{g\in G : g(\widetilde b)=\widetilde b\}$
its stabilizer. Since $G$ is abelian, this subgroup does not depend on the
choice of $\widetilde b$. Equivalently, $I_b$ is the inertia
group of the cover $\tC\to B$ at $b$. Since $C_{\iota\sigma}/B$ is unramified at $b$, we have $I_b\subseteq\langle\iota\sigma\rangle.$
On the other hand, since $C_{\sigma}/B$ is ramified at $b$, we have $I_b\not\subseteq\langle\sigma\rangle.$
The cyclic group $\langle\iota\sigma\rangle$ has order $4$, and its unique nontrivial proper subgroup is $\langle\sigma^2\rangle\subseteq\langle\sigma\rangle$. Hence we must have $I_b=\langle\iota\sigma\rangle\cong\ZZ/4\ZZ.$ Thus, there are at least $2d-2$ branch points of $\tC/B$ with this inertia group.

Now set
$$ C_{\iota\sigma^2}:=\tC/\langle\iota\sigma^2\rangle, \ \ \ C_{\iota,\sigma^2}=\tC/\langle\iota,\sigma^2\rangle.$$
Then $C_{\iota\sigma^2}\to C_{\iota,\sigma^2}$ is a double cover. We claim that $g(C_{\iota\sigma^2})-g(C_{\iota,\s^2})=1.$ Indeed, on the $\iota$-invariant differentials the invariance conditions for $\langle\iota\sigma^2\rangle$ and $\langle\iota,\sigma^2\rangle$ coincide. On the Prym part, invariance under $\iota\sigma^2$ is equivalent to $(\sigma^2)^*\omega=-\omega$. Since $\rho(\sigma)$ is a pseudoreflection, this is a one-dimensional condition. On the other hand, no Prym differential is invariant under $\iota$, so the claim is established.

Each point $b$ satisfying $I_b = \langle \iota\s \rangle$ gives a ramification point of $C_{\iota\sigma^2}\to C_{\iota, \sigma^2}$. Indeed, the fiber over $b$ of the cover $\tC \to B$ consists of two points $p_1$ and $p_2$ which are glued under $\tC \to C_{\iota\s^2}$ because $\langle \iota\s \rangle \cap \langle \iota\s^2\rangle = \{1\}$. Since the fiber of $C_{\iota\s^2} \to B$ consists of a unique point, this point must be a branch point of $C_{\iota\s^2} \to C_{\iota,\s^2}$.
Therefore, if $r$ is the ramification degree of $C_{\iota\s^2}\to C_{\iota,\s^2}$, then $r \geq 2d-2$.

On the other hand, by the Riemann-Hurwitz formula, we have $$r=4-2g(C_{\iota, \s^2})\leq4.$$
But $d\geq4$, so we get a contradiction with $r \geq 2d-2$. Thus, a lift of order $4$ is impossible.

It remains to rule out a lift of order $8$. Suppose that $\sigma$ has order $8$ and set $\tau:=\sigma^2$. Then $\tau$ has order $4$, while $\rho(\tau)=\alpha^2$
is a pseudoreflection of order 2. Moreover, $\tau^2=\sigma^4\neq1$
belongs to $\ker\rho$. By Lemma \ref{lem:kernel},
$$g\bigl(\tC/\langle\iota,\tau^2\rangle\bigr)\leq1.$$ Moreover, $\tau^2\neq\iota$, since
$\rho(\tau^2)=1$, whereas $\rho(\iota)=-\id_P$. Hence $\tau$ induces an automorphism of order 4 on $C$. Thus, $\tau$ satisfies the hypotheses of \cite[Lemma 3.5]{ALN}, which implies $g\leq3$. This contradicts $g\geq5$.

It remains to consider \(g=4\). Then
$d=\dim P=3$ and $g(\widetilde C)=7.$  Let
$\alpha\in\Aut(P,\Xi)$ be an order-$4$ pseudoreflection of geometric origin, and let $\sigma\in Z(\iota)$ be a lift of $\alpha$. We first observe that $ |\ker\rho|\leq3.$ Indeed, if $h=g(\widetilde C/\ker\rho)$, then, as in the proof of
\cite[Lemma 3.1]{ALN}, the fact that $\ker\rho$ acts trivially on $P$
implies $h\geq\dim P=3$. Hence the  Riemann-Hurwitz formula gives
$$
12=2g(\tC)-2\geq |\ker\rho|(2h-2)\geq4|\ker\rho|,$$ and the claim follows.

Therefore, $\sigma$ has order 4, 8 or 12. If $\sigma$ has order 12, then $\sigma^3$ has order 4, while
$\rho(\sigma^3)=\alpha^3$ is again an order-4 pseudoreflection. Thus, it suffices to exclude lifts of orders 4 and 8.

Suppose first that $\sigma$ has order 4. As in the proof above,

$$G:=\langle\iota,\sigma\rangle
\cong\mathbb{Z}/2\mathbb{Z}\times\mathbb {Z}/4\mathbb{Z}.$$
Set $B=\tC/G$. The same computation gives
$$g(C_\sigma)=g(B)+2,\ \ \ g(C_{\iota\sigma})=g(B),$$
and therefore
$r_\sigma=2(3-g(B))$ and $ r_{\iota\sigma}=2(1-g(B)).$
In particular, there are at least 4
branch points of $\tC\to B$ whose inertia group is $\langle\iota\sigma\rangle$ of order 4.

We claim that $g(B)=0$. Indeed, $r_{\iota\sigma}\geq0$, so
$g(B)\leq 1$. If $g(B)=1$, the four branch points above already
contribute $$4\cdot |G|\left(1-\frac14\right)=24$$
to the Riemann-Hurwitz formula for $\tC\to B$, while $2g(\widetilde C)-2=12$, which is a contradiction. Thus $B\simeq\mathbb{P}^1$.

Applying the Riemann-Hurwitz formula to 
$\tC\to\mathbb P^1$, we obtain $$12=-16+\sum_{b\in\mathbb P^1}
8\left(1-\frac1{|I_b|}\right).$$
Hence the total ramification contribution is 28. The four points
with inertia group $\langle\iota\sigma\rangle$ contribute 24. It follows that there are exactly four such points and exactly
one additional branch point, whose inertia group has order 2.
Thus the $G$-cover has signature $(0;4,4,4,4,2).$ 

Consider the local monodromies of this cover. The first four lie in
the cyclic subgroup $H:=\langle\iota\sigma\rangle\subsetneq G.$
Since the product of the five local monodromies is the identity, the
fifth monodromy also lies in $H$. Hence all local monodromies belong
to $H$. On the other hand, since the cover is connected and the base
is $\mathbb P^1$, its local monodromies must generate $G$. This is
a contradiction.

It remains to exclude the case in which $\sigma$ has order 8. Let $\tau=\sigma^2$. Then $\tau$ has order $4$, $\rho(\tau)=\alpha^2$ is a
pseudoreflection of order 2, and $\mu:=\tau^2=\sigma^4$ is a nontrivial involution acting trivially on $P$. We claim that $g\bigl(\tC/\langle\iota,\mu\rangle\bigr)\leq1.$ Indeed, the Prym variety of the double cover $\tC/\langle\iota\mu\rangle \to \tC/\langle\iota,\mu\rangle$
corresponds to the subspace of Prym differentials on which $\mu$
acts as $-1$. Since $\mu\in\ker\rho$, this subspace is zero. Thus, by the Riemann-Hurwitz formula, $g(\tC/\langle\iota,\mu\rangle) \leq 1$. Moreover, $\tau^2\neq\iota$, since
$\rho(\tau^2)=1$, whereas $\rho(\iota)=-\id_P$. Hence $\tau$ induces an automorphism of order 4 on $C$ so it satisfies all the hypotheses of \cite[Lemma 3.5]{ALN}. That lemma gives $g\leq3$, contradicting $g=4$.

Therefore an order-4 pseudoreflection cannot occur for $g=4$ either, and the result follows.
\end{proof}

\begin{exam}
\label{ex:order-four}
The bound in Theorem \ref{thm:order-two} is sharp. Indeed, there exist étale double covers of curves of genus 3 whose Prym varieties admit a pseudoreflection of geometric origin of order 4.

Let $$\widetilde G=\langle \iota,\sigma\rangle
\cong \mathbb Z/2\mathbb Z\times\mathbb Z/4\mathbb Z,$$
where $\iota$ has order 2 and $\sigma$ has order 4. Consider a connected $\widetilde G$-Galois cover $\tC\longrightarrow \mathbb{P}^1$ with local monodromies $$\iota\sigma,\ \ \ \iota\sigma,\ \ \ \sigma^2,\ \ \ \iota\sigma^2,\ \ \ \iota\sigma^2.$$
Such a cover exists by the Riemann existence theorem: these elements generate $\widetilde G$ and their product is the identity. Their orders are $(4,4,2,2,2)$, respectively. Hence the Riemann-Hurwitz formula gives
$$2g(\tC)-2 = -16+2\cdot 8\left(1-\frac14\right)+3\cdot 8\left(1-\frac12\right)=8,$$ and hence $g(\tC)=5$.

None of the inertia groups contains $\iota$, so $\iota$ acts freely on $\tC$. Thus $C:=\tC/\langle\iota\rangle$
has genus 3, and $\widetilde C\to C$ is an étale double cover. Let $P=P(\widetilde C/C)$. We claim that the automorphism of $P$ induced by $\sigma$ is a pseudoreflection of order $4$. Set $B=\tC/\widetilde G\simeq\mathbb P^1$. The quotient $\tC/\langle\sigma\rangle\longrightarrow B$
is a double cover branched at four points, and hence $g\bigl(\tC/\langle\sigma\rangle\bigr)=1$. On the other hand, $\tC/\langle\iota\sigma\rangle\to B$
is a double cover branched at two points, so
$g\bigl(\tC/\langle\iota\sigma\rangle\bigr)=0$. Let $$V=H^0(P,\Omega_P^1)=H^0(\widetilde C,\omega_{\tC})^{-}_{\iota}.$$
Since $g(B)=0$, we have $\dim V^{\sigma}= g(\tC/\langle \sigma \rangle) - g(B) =1.$
Moreover, on $V$ the involution $\iota$ acts as $-1$, and hence the $(-1)$-eigenspace of $\sigma$ is precisely the subspace on which $\iota\sigma$ acts trivially. Therefore, $$\dim V^{-}_{\sigma}=g\bigl(\tC/\langle\iota\sigma\rangle\bigr)-g(B)=0.$$
Since $\dim V=2$ and $\sigma^4=1$, the eigenvalues of $\sigma$ on $V$ are $1$ and $\pm i$, hence it is a pseudoreflection of order 4.
\end{exam}


\section{The faithful smooth-quotient classification}\label{sec:groups}

We recall from [ALA22, Theorem 3.5] that if \((P, \Xi)\) is a principally polarized abelian variety and \(H \leq \operatorname{Aut}(P, \Xi)\) is such that \(P / H\) is smooth, then there is an isogeny \(P \simeq\) \(X \times Y\) where \(H\) acts trivially on \(X\), and on \(Y\) with finitely many fixed points. Moreover, there exist elliptic curves \(E_1, \ldots, E_r, F_1, \ldots, F_s\) such that
\[
Y \cong \prod_{i=1}^r E_i^{n_i} \times \prod_{j=1}^s F_j^{t_j}
\]

\[
H=\prod_{i=1}^r\bigl(\left(\mathbb{Z} / m_i \mathbb{Z}\right)^{n_i} \rtimes S_{n_i}\bigr) \times \prod_{j=1}^s S_{t_j+1}
\]
where \(m_i \in\{2,3,4,6\}, n_i, t_j \geq 1\), and each factor of \(H\) acts on the corresponding factor of \(Y\) in an explicit way.

We now prove Theorem \ref{thm:main}, the main result of the paper.


\begin{thm}
\label{thm41}
Let $g\geq5$ and let $G\leq \Aut(\widetilde C)$ be a finite group such that every element of $G$ commutes with $\iota$. Let
$$\rho: G\longrightarrow \Aut(P,\Xi)$$
be the induced representation. Assume that $\rho$ is injective and that $P/\rho(G)$ is smooth. If $G$ is nontrivial, then $$G\cong \ZZ/2\ZZ \ \ \ \text{or}\ \ \ G\cong (\ZZ/2\ZZ)^2.$$
Moreover, if $g > 7$ then $G\cong \ZZ/2\ZZ$. 
Both groups occur for some $g \geq 5$.
\end{thm}

\begin{proof}
In the proof we identify $G$ with $\rho(G)$. According to \cite[Proposition 5.1]{ALN}, there are nine possible groups: $$\ZZ/2\ZZ, (\ZZ/2\ZZ)^2, \ZZ/4\ZZ, (\ZZ/4\ZZ)^2, \ZZ/2\ZZ \times \ZZ/4\ZZ, (\ZZ/2\ZZ)^3, (\ZZ/2\ZZ)^2\rtimes S_2, (\ZZ/4\ZZ)^2\rtimes S_2, \bigl((\ZZ/2\ZZ)^2\rtimes S_2\bigr)\times S_2.$$

We eliminate the seven groups other than $\ZZ/2\ZZ$ and $(\ZZ/2\ZZ)^2$.

\medskip
\noindent\textit{Step 1: Groups involving $4$-torsion.}

Since $P/G$ is smooth, $G$ is generated by pseudoreflections. By Theorem \ref{thm31}, all of these generators are involutions. This eliminates $\ZZ/4\ZZ, (\ZZ/4\ZZ)^2, \ZZ/2\ZZ\times \ZZ/4\ZZ$, since these groups are not generated by involutions. The only other possible group containing $4$-torsion is 
$G=(\ZZ/4\ZZ)^2\rtimes S_2,$ where $S_2$ exchanges the two $\ZZ/4$ factors. The homomorphism
$$\varphi\colon G\longrightarrow\ZZ/4\ZZ, \ \ \  \varphi(a,b;\eps)=a+b$$
is surjective. If $x\in G$ is an involution, then
\[
2\varphi(x)=\varphi(x^2)=0,
\]
so $\varphi(x)\in\{0,2\}$. Hence the subgroup generated by all involutions has image contained in the proper subgroup $\{0,2\}\subset\ZZ/4\ZZ$. Thus, $G$ is not generated by involutions. We have eliminated all four candidates involving $\ZZ/4\ZZ$.

\medskip
\noindent\textit{Step 2: The dihedral cases.}

Suppose first that $G\cong (\ZZ/2\ZZ)^2\rtimes S_2$ is the dihedral group of order 8. The group $G$ is not a direct product of two groups and $|G| = 8$, hence by \cite[Theorem 3.5]{ALA}, we have $$(P,\Xi)\cong(E^2,\Theta_2)\times(B,\Theta_B)$$ as principally polarized abelian varieties. Here, $E$ is an elliptic curve, with $G$ acting by the two coordinate sign changes and by permutation of the two factors on $E^2$ and trivially on $B$. Since $d\geq4$, we have $\dim B=d-2>0$, therefore $\Xi$ has at least three irreducible components.

Moreover, since $P$ is decomposable, it follows from \cite[Remark 4.2]{ALN} that
$$(P,\Xi)\cong(F,[0])\times(JH,\Theta_H)$$ for an elliptic curve $F$ and a hyperelliptic curve $H$. But $\Theta_H$ is irreducible, hence $\Xi$ has two irreducible components, which is a contradiction. Hence $(\ZZ/2\ZZ)^2\rtimes S_2$ cannot occur. The same argument excludes $\bigl((\ZZ/2\ZZ)^2\rtimes S_2\bigr)\times S_2$.

\medskip
\noindent\textit{Step 3: The group $(\ZZ/2\ZZ)^3$.}

It remains to rule out $G \cong (\ZZ/2\ZZ)^3$. Let $\sigma_1, \sigma_2, \sigma_3$ be pseudoreflections generating $G$. Then we have $$H^0(P,\Omega_P^1)=W\oplus L_1\oplus L_2\oplus L_3,$$
where $\dim W=d-3, \dim L_i=1$ and $\sigma_i$ acts by $-1$ on $L_i$ and by $+1$ on the other summands. Since $d\geq4$, the common fixed space $W$ is nonzero. Every element of $G$ acts trivially on $W$, whereas $\iota$ acts as $-1$ on all Prym differentials, hence $\iota\notin G$. Let $$K:=\langle\iota,G\rangle\cong(\ZZ/2\ZZ)^4.$$
Set $B:=\widetilde C/K,$ and consider the eight characters $$\chi: K\longrightarrow\{\pm1\}$$
with $\chi(\iota)=-1$. For each such character, let $C_\chi:=\tC/\ker\chi$. If $V_\chi$ denotes the $\chi$-eigenspace in $H^0(\tC,K_{\tC})$, then $$ H^0(C_\chi,K_{C_\chi})=H^0(B,K_B)\oplus V_\chi,$$
so $m_{\chi} := \dim V_\chi=g(C_\chi)-g(B).$

Let $n_\chi$ be the number of branch points of $C_\chi\to B$. The Riemann-Hurwitz formula gives $$n_\chi=2(m_\chi-g(B)+1).$$ Let $R$ be the number of branch points of the full $K$-cover $\widetilde C\to B$. Recall that for the action of a group on a smooth curve all inertia subgroups are cyclic, thus, since $K$ is elementary abelian, every nontrivial inertia subgroup has order $2$. Its generator is never $\iota$, because the covering involution is fixed-point-free. For any nontrivial inertia element $v\neq\iota$, exactly four of the eight characters with $\chi(\iota)=-1$ satisfy $\chi(v)=-1$. Thus, every branch point of $\tC/B$ is counted in exactly four of the double covers $C_\chi/B$, and
$$\sum_\chi n_\chi=4R.$$
Thus, we obtain $$R=\frac{d+8-8g(B)}{2}.$$ In particular, $d$ is even.

Now take the character $\chi_0$, which is trivial on $G$, i.e. corresponding to the common fixed space $W$. Its eigenspace has dimension $d-3$, and therefore $n_{\chi_0}=2(d-g(B)-2).$
Every branch point of the corresponding double cover is a branch point of $\tC/B$, so $n_{\chi_0}\leq R$. Therefore,

$$2(d-g(B)-2)\leq\frac{d+8-8g(B)}{2},$$
which simplifies to

$$3d+4g(B)\leq16.$$
Since $d\geq4$ is even, we must have $d = 4$.

By \cite[Theorem 3.5 and Lemma 3.6]{ALA}, each of the three $\ZZ/2\ZZ$-factors gives either a principally
polarized elliptic direct factor, or a one-dimensional factor of the standard
construction, in which case the restricted polarization is of type $(2)$.
Let $p$ and $q$ denote the numbers of factors of these two types,
respectively. 

Again, by \cite[Remark 4.2]{ALN}, the only case in which a Prym variety admitting
a geometric pseudoreflection is a polarized product is $$(P,\Xi)\simeq (E,\Theta_E)\times(JH,\Theta_H).$$
Therefore, $p\leq 1$, and thus $q\geq2$.

For $q$ one-dimensional standard factors, the moving part has dimension
$q$ and polarization type $(2,\ldots,2)$. In the standard construction from \cite{ALA}, the auxiliary polarized abelian variety has dimension at least $q$. Consequently, $$d\ge p+2q.$$
Since $p+q=3$ and $p\leq 1$, we obtain
$d \geq 5$
contradicting \(d=4\). Thus \(G\simeq(\ZZ/2\ZZ)^3\) cannot occur.

\medskip
\noindent\textit{Step 4: $G \cong (\mathbb{Z}/2\mathbb{Z})^2$ forces $g \leq 7$.}

Let $G \cong (\mathbb{Z}/2\mathbb{Z})^2$ and $B := \tC / (G \times \langle\iota\rangle)$. Repeating the analysis from the first part of Step 3, we get the inequality $$2(g - g(B) - 2) \leq d - 4g(B) + 4,$$ which is equivalent to $$d+2g(B) \leq 6.$$ Thus, $d = g-1 \leq 6$, hence $g \leq 7$.

Finally, we note that both groups $\ZZ/2\ZZ$ and $(\ZZ/2\ZZ)^2$ occur by the constructions in \cite{ALN}; see \cite[Remark 5.2]{ALN} and also \cite[Remark 2.7]{BO26}. This proves the theorem.
\end{proof}


\begin{thebibliography}{99}
\bibitem[ALA]{ALA} R.~Auffarth and G.~Lucchini Arteche, \emph{Smooth quotients of principally polarized abelian varieties}, Mosc. Math. J. \textbf{22} (2022), no.~2, 225--237.
\bibitem[ALN]{ALN} R.~Auffarth, M.~Lahoz and J.~C.~Naranjo, \emph{Pseudoreflections on Prym varieties}, arXiv:2412.04940v2.
\bibitem[BO26]{BO26} P.~Borówka and A.~Ortega, \emph{Klein coverings over hyperelliptic genus 3 curves}, arXiv:2602.18969, 2026.
\end{thebibliography}
\end{document}
